\scanpage{2}

% printed form
3 3 -- N O D E E 3 3

S E333

\handwritten{\[
\overset{\text{$Z_a$}}{\partial\longrightarrow I}
\]}

\handwritten{$\deleted{\unsure{\mathrm{ob}}}\,\unsure{\pi_0|\Pi_0}(\partial)\geq 2$\quad \deleted{\unsure{0-connexe|$\varnothing$-connexe}} \qquad \unsure{0-connexe}}

\handwritten{\[
\left\{\begin{aligned}
&\phi\longrightarrow e\\
&\partial\times S^n(N)\longrightarrow I\times S^n(N)\\
&\bigl\{+\,\text{Cônes gauches (resp. droits) sur les précédents}\\
&\text{(donne aussi }\{0\}\to\Delta^1\\
&\qquad\text{resp.}^{\mathrm{t}}\ \{1\}\to\Delta^1\text{)}
\end{aligned}\right.
\]}

\handwritten{\[
\underline{\mathrm{Hom}}(\unsure{\pi|\Pi|T|I},X)
\]}

\handwritten{\uline{Lemme} a) Si $X\longrightarrow e$ a la propriété RLP pour $\phi\longrightarrow e$ \uline{et} les $\partial\times S^n(N)\longrightarrow I\times S^n(N)$, $n\in\mathscr N$ (\unsure{partie} infinie de $N$), alors $X$ est $W_\infty$-asph.}

\handwritten{b) Si $X\longrightarrow Y$ a la RLP p.r. aux cônes gauches (resp. droits) p.r. aux morphismes précédents, alors $X$ est $W_\infty$-propre (resp. $W_\infty$-lisse).}

\handwritten{\uline{Cor} a) Si \inserted{$X\longrightarrow Y$} a la propriété RLP p.r. aux \unsure{sommes} de deux \deleted{\unsure{autres}} \inserted{(resp. \unsure{jointures|gauches})} \deleted{\unsure{entre}} catégories \deleted{\unsure{fibrées}} \uline{réunies}, \uline{finies}, et alors c'est un morphisme $W_\infty$-propre (resp. lisse) \inserted{à fibres $W_\infty$-asphériques}.}

\handwritten{b) Si $f\colon X\longrightarrow Y$ a la propriété RLP p.r. aux \inserted{\unsure{demi-cercles} (resp. \unsure{demi}) $n$\textsuperscript{es} \deleted{\unsure{triviales}} \unsure{ét.c.}} \unsure{sommes} de deux (qui s.t des $W_\infty$, \inserted{des cat finies}) alors $f$ est $W_\infty$-propre (resp. $W_\infty$-lisse). S'il a la dite propriété pour les $W_\infty$ cofibrations triviales \inserted{entre cat. finies}, alors il est $W_\infty$-parfait.}

% archive sheet number: 1
